\subsubsection{Random SPIN exponent closure}
\label{app:random-spin-exponents}

This subsection proves the two estimates used by
Theorem~\ref{thm:independent-random-spin}. The outer exponent follows from
the exact injection law in~\eqref{eq:injection-global-mass}. For a sequence
$h/N\to\eta\in(0,1)$, coefficient extraction gives
\begin{equation}
  \Phi_{\mathrm{blk}}(\eta)
  =
  \begin{cases}
    \eta\log_2(1+\sqrt2),
      &0<\eta\le1-2^{-1/2},\\[1mm]
    H_2(\eta)-1/2,
      &1-2^{-1/2}\le\eta<1.
  \end{cases}
  \label{eq:random-outer-linear-exponent}
\end{equation}

The inner exponent comes from the transfer matrix in
\eqref{eq:random-transfer-matrix}. For $a>0$ and $S\in(0,1)$, define
\[
  \lambda_\infty(a,S)
  :=
  \max\left\{1,\frac{(1+a)(1+S)}2\right\}.
\]

\begin{lemma}[Limiting convolution growth]
\label{lem:limiting-convolution-growth}
If $m_N=\Theta(\log N)$, then
\begin{equation}
  \limsup_{N\to\infty}\frac1N\log_2
  \left(
    e_{m_N}^{\mathsf T}T_{m_N}(a,S)^N\boldsymbol 1
  \right)
  \le\log_2\lambda_\infty(a,S).
  \label{eq:limiting-convolution-growth}
\end{equation}
\end{lemma}

\begin{proof}
Set $b:=(1+a)/2$ and let $\lambda_m$ be the Perron root of
$T_m(a,S)$. The right Perron-vector equations give
\begin{equation}
  1
  =
  \frac{bS}{\lambda_m}
  \frac{1-(b/\lambda_m)^m}{1-b/\lambda_m}
  +
  \left(\frac b{\lambda_m}\right)^m
  \frac{aS}{\lambda_m-1}.
  \label{eq:perron-root-identity}
\end{equation}
If $b(1+S)<1$, this identity gives $\lambda_m\to1$. If
$b(1+S)>1$, the active-state principal submatrix and
\eqref{eq:perron-root-identity} give
$\lambda_m\to b(1+S)$. At equality, any subsequential limit greater than one
contradicts~\eqref{eq:perron-root-identity}.

Normalize the Perron vector at state zero. The same equations show that the
ratio of its largest and smallest coordinates is $\exp(O(m))$. Comparing the
all-ones vector with coordinatewise multiples of this Perron vector bounds the
matrix power by $\lambda_m^N\exp(O(m))$. Since $m_N=O(\log N)$, this factor is
subexponential and proves~\eqref{eq:limiting-convolution-growth}.
\end{proof}

For $\delta\in(0,1/2)$, equations
\eqref{eq:random-convolution-tail} and
\eqref{eq:limiting-convolution-growth} give the inner exponent
\begin{equation}
  F_\delta(\eta)
  :=
  \inf_{\substack{a>0\\0<S<1}}
  \left[
    \log_2\lambda_\infty(a,S)
    -\eta\log_2a
    -\delta\log_2S
    -H_2(\eta)
  \right].
  \label{eq:random-inner-linear-exponent}
\end{equation}
Set
\[
  \eta_{\mathrm{conv}}(\delta)
  :=1-\frac1{2(1-\delta)},
  \quad
  a_\eta
  :=\frac{\eta}{\sqrt{\eta^2+\delta^2}+\delta},
  \quad
  S_\eta:=\frac{1-a_\eta}{1+a_\eta}.
\]
Direct minimization yields
\begin{equation}
  F_\delta(\eta)
  =
  \begin{cases}
    -\eta\log_2a_\eta-\delta\log_2S_\eta-H_2(\eta),
      &0<\eta\le\eta_{\mathrm{conv}}(\delta),\\[1mm]
    H_2(\delta)-1,
      &\eta_{\mathrm{conv}}(\delta)\le\eta<1.
  \end{cases}
  \label{eq:closed-random-inner-exponent}
\end{equation}

\begin{lemma}[Linear-weight closure]
\label{lem:random-spin-linear-closure}
Set $\delta_*:=5501/50000$ and
$J_{\delta_*}:=\Phi_{\mathrm{blk}}+F_{\delta_*}$. Then
\[
  J_{\delta_*}(\eta)<0
  \qquad(0<\eta<1).
\]
The inequality is uniform on every compact subinterval of $(0,1)$.
\end{lemma}

\begin{proof}
Set $\eta_{\mathrm{blk}}:=1-2^{-1/2}$. On
$(0,\eta_{\mathrm{blk}}]$, the derivative of $J_{\delta_*}$ is increasing,
so $J_{\delta_*}$ is convex and $J_{\delta_*}(0)=0$. Outward interval
arithmetic gives
\[
  J_{\delta_*}(\eta_{\mathrm{blk}})
  \in[-0.074030,-0.074029].
\]
Convexity proves strict negativity on the first region. On
$[\eta_{\mathrm{blk}},\eta_{\mathrm{conv}}(\delta_*)]$, the entropy terms
cancel and the remaining expression is increasing. Its right endpoint lies
in $[-0.011076,-0.011075]$. On the final region,
\[
  J_{\delta_*}(\eta)
  =H_2(\eta)+H_2(\delta_*)-\frac32
  \le H_2(\delta_*)-\frac12,
\]
and the last value lies in
$[-0.000023719,-0.000023718]$. These strict intervals prove the claim.
\end{proof}

The remaining issue is uniformity as $\eta\downarrow0$. Define
\[
  \eta_*:=\frac1{5000},
  \qquad
  \chi:=\frac{101}{100},
  \qquad
  \gamma:=\frac{51}{50}.
\]
Set
\begin{align*}
  a_*&:=\frac{\eta_*}{\delta_*(\chi+1)},
  &b_*&:=\frac{1+a_*}{2},\\
  \epsilon_*&:=
  \frac{\chi^2a_*}{(\chi+1)(1-\chi a_*)}
  +\frac{\eta_*}{2(1-\eta_*)},
  &\rho_*&:=\delta_*(\chi+1)e^{\epsilon_*}.
\end{align*}

\begin{lemma}[Sparse-weight closure]
\label{lem:random-spin-sparse-closure}
Let $m_N:=\lceil\gamma\log_2N\rceil$. There are constants $C,N_0$ such
that, for $N\ge N_0$ and $1\le h\le\eta_*N$,
\begin{equation}
  p_{N,h}(\delta_*)
  \le C N^{2\gamma}\sqrt h\,\rho_*^h.
  \label{eq:random-spin-sparse-tail}
\end{equation}
If $B_N$ is even, divides $N$, and satisfies
$B_N\ge17\log_2N$, then
\begin{equation}
  \sum_{h=1}^{\lfloor\eta_*N\rfloor}
  \E[A_{N,h}^{\mathrm{out}}]p_{N,h}(\delta_*)
  =o(1).
  \label{eq:random-spin-sparse-sum}
\end{equation}
\end{lemma}

\begin{proof}
For $\eta:=h/N$, evaluate~\eqref{eq:random-convolution-tail} at
\[
  a:=\frac{\eta}{\delta_*(\chi+1)},
  \qquad
  S:=\frac{1-\chi a}{1+a}.
\]
The active-state row sum equals $1-(\chi-1)a/2<1$. The Perron-root identity
then bounds the matrix power by $CN^{2\gamma}$ uniformly over the stated
range. A standard lower bound on $\binom Nh$ gives
\eqref{eq:random-spin-sparse-tail}; the exponent correction increases with
$\eta$ and is at most $\epsilon_*$. Interval arithmetic certifies
\begin{align*}
  \rho_*&\in
  [0.2212639483446276,0.2212639483446278],\\
  \gamma\log_2(1/b_*)-1&>0.0186697.
\end{align*}
In particular, $\rho_*<\sqrt2-1$.

The exact independent-block outer law gives
\[
  \sum_{h\le\eta_*N}
  \E[A_{N,h}^{\mathrm{out}}]p_{N,h}(\delta_*)
  \le
  C N^{2\gamma+1/2}
  \left[(1+G_{B_N,1/2}(\rho_*))^{N/B_N}-1\right].
\]
Using~\eqref{eq:injection-local-mass}, the right side is bounded by a
constant times
\[
  \frac{N^{2\gamma+3/2}}{B_N}
  \left(\frac{1+\rho_*}{\sqrt2}\right)^{B_N}.
\]
Finally, outward arithmetic gives
\[
  17\log_2\!\left(\frac{\sqrt2}{1+\rho_*}\right)
  -(2\gamma+3/2)>0.0576243.
\]
Thus the last expression tends to zero.
\end{proof}

Appendix~\ref{app:structured-scalable-certificate} gives the structured-route
and IMT transfer derivations and arithmetic semantics.
The theorem uses the probability space and inequalities stated in the main
text and appendices. Our author-side verification records retain the verifier
sources and hash-bound arithmetic witnesses; these are outside the core
encoder artifact.
